\section{The Diagram Grammar Family}
\label{sec:diagram-family}

Building on the corpus taxonomy (Section~\ref{sec:corpus-taxonomy}), this section specifies the planned diagram grammars beyond Timeline. Each grammar has a distinct Domain IR, layout algorithm, and use case.

\subsection{Flow/Pipeline Grammar}
\label{sec:flow-grammar-overview}

The Flow grammar describes ordered sequences of steps connected by directional edges. It is the most natural extension of the timeline grammar. \textbf{The full concrete specification of the Flow Domain IR, layout contract, and lowering semantics is given in Section~\ref{sec:flow-grammar}.} This subsection provides the high-level sketch; the authoritative definition lives in the dedicated section.

\subsubsection{Domain IR Sketch}

\begin{lstlisting}[language=yaml,caption={Flow/Pipeline Domain IR (sketch)}]
version: "1.0"
metadata:
  title: "CI/CD Pipeline"
  theme: dark-flow
flow:
  direction: left-to-right | top-to-bottom
  nodes:
    - id: trigger
      label: "Push to main"
      icon: git-branch
      category: trigger
    - id: build
      label: "Build"
      icon: package
      category: step
    - id: test
      label: "Test"
      icon: check-circle
      category: step
    - id: deploy
      label: "Deploy"
      icon: rocket
      category: output
  edges:
    - from: trigger
      to: build
      style: solid
    - from: build
      to: test
      style: solid
      animation: flowing-dashes    # optional
    - from: test
      to: deploy
      style: solid
      label: "on success"
  lanes:                           # optional swimlanes
    - id: dev
      label: "Development"
      nodes: [trigger, build, test]
    - id: prod
      label: "Production"
      nodes: [deploy]
\end{lstlisting}

\subsubsection{Layout Algorithm}

\begin{enumerate}
  \item Parse nodes and edges into a directed graph
  \item If lanes are specified, partition nodes by lane
  \item Compute node positions: left-to-right (x by topological order, y by lane) or top-to-bottom
  \item Route edges between node ports (avoid crossing where possible)
  \item Apply animation hints to edges
\end{enumerate}

\subsubsection{Reuse from Timeline}

\begin{itemize}
  \item Icon rendering (icon registry)
  \item Node shapes (rounded rect, circle)
  \item Connector routing (straight, orthogonal, curved)
  \item Swimlane layout (horizontal tracks)
  \item Theme styling (colors, typography)
\end{itemize}

\subsection{Comparison Grammar}
\label{sec:comparison-grammar}

The Comparison grammar describes side-by-side evaluations of items across features.

\subsubsection{Domain IR Sketch}

\begin{lstlisting}[language=yaml,caption={Comparison Domain IR (sketch)}]
version: "1.0"
metadata:
  title: "Framework Comparison"
  theme: consulting
comparison:
  columns:
    - id: react
      label: "React"
      icon: react-logo
    - id: vue
      label: "Vue"
      icon: vue-logo
  rows:
    - id: performance
      label: "Performance"
      cells:
        react: { indicator: high, note: "Virtual DOM" }
        vue: { indicator: high, note: "Reactive system" }
    - id: learning-curve
      label: "Learning Curve"
      cells:
        react: { indicator: medium }
        vue: { indicator: low }
    - id: ecosystem
      label: "Ecosystem"
      cells:
        react: { indicator: high, note: "Largest" }
        vue: { indicator: medium }
  indicators:
    high: { icon: check, color: green }
    medium: { icon: minus, color: amber }
    low: { icon: x, color: red }
\end{lstlisting}

\subsubsection{Layout Algorithm}

\begin{enumerate}
  \item Compute column widths (equal or content-based)
  \item Compute row heights (based on cell content)
  \item Position column headers
  \item Position row labels and cells in grid
  \item Render indicators and notes within cells
\end{enumerate}

\subsection{Stat Callout Grammar}
\label{sec:stat-grammar}

The Stat grammar describes large numbers with captions, used for impact communication.

\subsubsection{Domain IR Sketch}

\begin{lstlisting}[language=yaml,caption={Stat Callout Domain IR (sketch)}]
version: "1.0"
stat:
  value: "98.7%"
  label: "Accuracy"
  sublabel: "PageIndex (Vectorless RAG)"
  trend:                          # optional
    direction: up
    delta: "+12.3%"
  comparison:                     # optional
    value: "31%"
    label: "GPT-4o Search"
\end{lstlisting}

\subsubsection{Layout Algorithm}

Trivial: center the large number, position label below, trend indicator to the side.

\subsection{Node-Link Graph Grammar}
\label{sec:graph-grammar}

The Graph grammar describes networks of vertices and edges. This is the most complex grammar due to layout challenges.

\subsubsection{Domain IR Sketch}

\begin{lstlisting}[language=yaml,caption={Node-Link Graph Domain IR (sketch)}]
version: "1.0"
metadata:
  title: "System Architecture"
  theme: architecture
graph:
  layout: tidy-tree | dag | hub-spoke | force
  vertices:
    - id: api
      label: "API Gateway"
      icon: server
      cluster: frontend
    - id: auth
      label: "Auth Service"
      icon: lock
      cluster: services
    - id: db
      label: "Database"
      icon: database
      cluster: data
  edges:
    - from: api
      to: auth
      style: solid
    - from: auth
      to: db
      style: dashed
  clusters:                       # optional grouping
    - id: frontend
      label: "Frontend"
    - id: services
      label: "Services"
    - id: data
      label: "Data Layer"
\end{lstlisting}

\subsubsection{Layout Algorithms}

Graph layout is the hard problem. The grammar supports multiple layout algorithms, each suited to different graph structures:

\begin{description}
  \item[tidy-tree] For trees (single root, no cycles). Uses Reingold-Tilford~\cite{reingold1981}.
  \item[dag] For directed acyclic graphs. Uses Sugiyama layered layout~\cite{sugiyama1981}.
  \item[hub-spoke] For star graphs (central hub with spokes). Custom radial layout.
  \item[force] For general graphs. Uses force-directed layout~\cite{fruchterman1991}. Note: force-directed is non-deterministic without careful seeding.
\end{description}

\paragraph{Determinism constraint.}
Force-directed layout involves iterative simulation that can converge to different local minima. To preserve determinism:
\begin{itemize}
  \item Initial positions derived from node IDs (sorted, then arranged on a circle)
  \item Random perturbations seeded by \texttt{scene\_hash}
  \item Fixed iteration count (no convergence-based termination)
\end{itemize}

\subsection{Step-Cards Grammar}
\label{sec:stepcards-grammar}

The Step-Cards grammar describes numbered instructional sequences.

\subsubsection{Domain IR Sketch}

\begin{lstlisting}[language=yaml,caption={Step-Cards Domain IR (sketch)}]
version: "1.0"
steps:
  layout: horizontal | vertical | grid
  cards:
    - ordinal: 1
      title: "Document Indexing"
      icon: file-text
      bullets:
        - "Parses the full document"
        - "No chunking required"
        - "No vector DB needed"
    - ordinal: 2
      title: "Tree-Based Reasoning"
      icon: git-branch
      bullets:
        - "LLM traverses tree"
        - "Evaluates context at each level"
\end{lstlisting}

\subsection{Grammar Interaction: Composition}

These grammars can be combined in a Composition IR (Section~\ref{sec:composition}). A single poster might contain:

\begin{itemize}
  \item A stat callout section (headline numbers)
  \item A comparison table (feature matrix)
  \item A flow diagram (process explanation)
  \item A timeline (roadmap)
\end{itemize}

Each section is a panel with its own grammar and Domain IR. The Composition layer arranges them.

\subsection{Prior Art and Tools}
\label{sec:diagram-prior-art}

The diagram grammar family builds on established prior art:

\begin{description}
  \item[Graphviz/DOT~\cite{graphviz}] The classic graph-description language. Excellent layout algorithms, poor visual quality for presentations.
  
  \item[Mermaid~\cite{mermaid2023}] Markdown-friendly diagramming. Good ecosystem integration, limited visual quality and theming.
  
  \item[PlantUML~\cite{plantuml}] UML-centric diagramming. Verbose syntax, technical aesthetic.
  
  \item[D2~\cite{d2lang}] Modern diagram scripting language with auto-layout. Closest in philosophy to this compiler.
  
  \item[Excalidraw] Hand-drawn aesthetic, manual placement. Not declarative.
  
  \item[Structurizr/C4~\cite{structurizr}] Domain-specific to software architecture.
\end{description}

The diagram compiler differentiates by combining: (1) presentation-quality output, (2) determinism, (3) theming, (4) agent-friendly schemas, and (5) multiple backends (SVG, PNG, PPTX).
